% GENERATED FILE. Edit canonical lessons, then run npm run generate:books.
% canonical-source-hash: fnv1a64:99a0a40a7f94556b
% canonical-lessons: KA-C203-belesu, KA-C203-aralu, KA-C203-badu, KA-C203-bisu, KA-C203-karagu

\chapter{To grow, To bloom, To wilt, To blow, To melt}
\label{ch:ka-to-grow-to-bloom-to-wilt-to-blow-to-melt}
\hlchaptermodality{203}

\begin{chapteropening}
\textbf{By the end of this chapter:} \emph{I can say to grow, to bloom, to wilt, to blow and to melt.}

Complete the last lesson of chapter 203: I can say to grow, to bloom, to wilt, to blow and to melt.
\end{chapteropening}

\section[\texorpdfstring{\kn{ಬೆಳೆಸು} (beḷesu)}{ಬೆಳೆಸು (beḷesu)}]{\kn{ಬೆಳೆಸು} (beḷesu) — to grow (plants), to raise}

\label{lesson:KA-C203-belesu}



\begin{quote}
Before the new one: say the Kannada for to contact, then the Kannada for to discuss.
\end{quote}

\subsection*{You'll want to know: \kn{ಬೆಳೆಸು}}
\textbf{\kn{ಬೆಳೆಸು}} — \emph{beḷesu} — \textquotedblleft{}to grow (plants), to raise\textquotedblright{}.

\begin{grammarlens}[title={What you've built}]
One more word for this chapter.
\end{grammarlens}

\subsection*{Your turn}
\begin{itemize}
\raggedright
  \item \emph{Say it:} \emph{beḷesu}
  \item \emph{Say it:} \emph{beḷesu}, once more
  \item \emph{Say it:} say \emph{beḷesu}
  \item \emph{Recall:} say the Kannada for to contact, then the Kannada for to discuss, then say \emph{beḷesu} again
\end{itemize}

\begin{tcolorbox}[breakable,colback=teal!4,colframe=teal!35!black,title={Before you move on}]
What does \kn{ಬೆಳೆಸು} mean? (\textquotedblleft{}To grow (plants), to raise\textquotedblright{}.) Say it once more.
\end{tcolorbox}
\section[\texorpdfstring{\kn{ಅರಳು} (araḷu)}{ಅರಳು (araḷu)}]{\kn{ಅರಳು} (araḷu) — to bloom}

\label{lesson:KA-C203-aralu}



\begin{quote}
Before the new one: say the Kannada for to discuss, then the Kannada for to grow.
\end{quote}

\subsection*{You'll want to know: \kn{ಅರಳು}}
\textbf{\kn{ಅರಳು}} — \emph{araḷu} — \textquotedblleft{}to bloom\textquotedblright{}.

\begin{grammarlens}[title={What you've built}]
One more word for this chapter.
\end{grammarlens}

\subsection*{Your turn}
\begin{itemize}
\raggedright
  \item \emph{Say it:} \emph{araḷu}
  \item \emph{Say it:} \emph{araḷu}, once more
  \item \emph{Say it:} say \emph{araḷu}
  \item \emph{Recall:} say the Kannada for to discuss, then the Kannada for to grow, then say \emph{araḷu} again
\end{itemize}

\begin{tcolorbox}[breakable,colback=teal!4,colframe=teal!35!black,title={Before you move on}]
What does \kn{ಅರಳು} mean? (\textquotedblleft{}To bloom\textquotedblright{}.) Say it once more.
\end{tcolorbox}
\section[\texorpdfstring{\kn{ಬಾಡು} (bāḍu)}{ಬಾಡು (bāḍu)}]{\kn{ಬಾಡು} (bāḍu) — to wilt, to wither}

\label{lesson:KA-C203-badu}



\begin{quote}
Before the new one: say the Kannada for to grow, then the Kannada for to bloom.
\end{quote}

\subsection*{You'll want to know: \kn{ಬಾಡು}}
\textbf{\kn{ಬಾಡು}} — \emph{bāḍu} — \textquotedblleft{}to wilt, to wither\textquotedblright{}.

\begin{grammarlens}[title={What you've built}]
One more word for this chapter.
\end{grammarlens}

\subsection*{Your turn}
\begin{itemize}
\raggedright
  \item \emph{Say it:} \emph{bāḍu}
  \item \emph{Say it:} \emph{bāḍu}, once more
  \item \emph{Say it:} say \emph{bāḍu}
  \item \emph{Recall:} say the Kannada for to grow, then the Kannada for to bloom, then say \emph{bāḍu} again
\end{itemize}

\begin{tcolorbox}[breakable,colback=teal!4,colframe=teal!35!black,title={Before you move on}]
What does \kn{ಬಾಡು} mean? (\textquotedblleft{}To wilt, to wither\textquotedblright{}.) Say it once more.
\end{tcolorbox}
\section[\texorpdfstring{\kn{ಬೀಸು} (bīsu)}{ಬೀಸು (bīsu)}]{\kn{ಬೀಸು} (bīsu) — to blow (of wind)}

\label{lesson:KA-C203-bisu}



\begin{quote}
Before the new one: say the Kannada for to bloom, then the Kannada for to wilt.
\end{quote}

\subsection*{You'll want to know: \kn{ಬೀಸು}}
\textbf{\kn{ಬೀಸು}} — \emph{bīsu} — \textquotedblleft{}to blow (of wind)\textquotedblright{}.

\begin{grammarlens}[title={What you've built}]
One more word for this chapter.
\end{grammarlens}

\subsection*{Your turn}
\begin{itemize}
\raggedright
  \item \emph{Say it:} \emph{bīsu}
  \item \emph{Say it:} \emph{bīsu}, once more
  \item \emph{Say it:} say \emph{bīsu}
  \item \emph{Recall:} say the Kannada for to bloom, then the Kannada for to wilt, then say \emph{bīsu} again
\end{itemize}

\begin{tcolorbox}[breakable,colback=teal!4,colframe=teal!35!black,title={Before you move on}]
What does \kn{ಬೀಸು} mean? (\textquotedblleft{}To blow (of wind)\textquotedblright{}.) Say it once more.
\end{tcolorbox}
\section[\texorpdfstring{\kn{ಕರಗು} (karagu)}{ಕರಗು (karagu)}]{\kn{ಕರಗು} (karagu) — to melt}

\label{lesson:KA-C203-karagu}



\begin{quote}
Before the new one: say the Kannada for to wilt, then the Kannada for to blow.
\end{quote}

\subsection*{You'll want to know: \kn{ಕರಗು}}
\textbf{\kn{ಕರಗು}} — \emph{karagu} — \textquotedblleft{}to melt\textquotedblright{}.

\begin{grammarlens}[title={What you've built}]
One more word for this chapter.
\end{grammarlens}

\subsection*{Your turn}
\begin{itemize}
\raggedright
  \item \emph{Say it:} \emph{karagu}
  \item \emph{Say it:} \emph{karagu}, once more
  \item \emph{Say it:} say \emph{karagu}
  \item \emph{Recall:} say the Kannada for to wilt, then the Kannada for to blow, then say \emph{karagu} again
\end{itemize}

\begin{tcolorbox}[breakable,colback=teal!4,colframe=teal!35!black,title={Before you move on}]
What does \kn{ಕರಗು} mean? (\textquotedblleft{}To melt\textquotedblright{}.) Say it once more.
\end{tcolorbox}
